% ===========================================================================
\section{Data and accounting boundaries}
\label{sec:data}
% ===========================================================================

\subsection{FY2015--16 to FY2023--24 infrastructure controls}

Oil and gas production come from the Pakistan Economic Survey 2025--26, which reports
\num{25.812} million barrels of crude oil and \num{1.141} million MMCF of natural gas in
FY2023--24 \citep{peseconomic2026}. Converted to IPCC activity units, these are
\num{4103.77} thousand m$^3$ of oil and
\num{32309.47} million m$^3$ of gas.

The Oil and Gas Regulatory Authority (OGRA) archive supplies annual refinery, network, customer,
consumption and UFG controls from FY2015--16 to FY2023--24
\citep{ograarchive,ogra2022,ogra2023,ogra2024}. Values are transcribed from the reported annual tables with
no interpolation. The image-based FY2020--21 report was independently OCR-audited against the
page images, and the FY2016--17 refinery total is retained as the regulator's restated value.
In FY2023--24, SNGPL reports
\SI{9320.69}{\kilo\metre} of transmission and
\SI{151419}{\kilo\metre} of distribution pipeline; SSGCL reports \SI{4206}{} and
\SI{49811}{\kilo\metre}, respectively. The combined systems serve \num{10748538} domestic,
\num{91925} commercial and \num{11318} industrial customers. Seventeen named compressor sites,
two LNG import terminals and \num{3437.13} MMCF d$^{-1}$ of utility and independent-system sales are
reported. The reference T\&D estimate uses measured pipeline length; throughput is an alternative
IPCC basis. Across the nine-year record, transmission length rises from
\SI{11555.08}{\kilo\metre} to \SI{13526.69}{\kilo\metre}, distribution length from
\SI{148278}{\kilo\metre} to \SI{201230}{\kilo\metre}, and domestic customers from \num{7997096} to
\num{10748538}. This supports a nine-year segment-resolved reconstruction rather than a
production-only historical extrapolation.

The FY2021--22 report identifies \SI{8200}{\kilo\metre} of polythene distribution pipe in the
SSGCL system. This is only \SI{4.2}{\percent} of the reported national distribution network and
does not classify the SNGPL network. Utility reports document active loss-control work---SNGPL
surveyed \SI{50250}{\kilo\metre} of underground network in FY2023--24 and both utilities report
rehabilitation and leak surveys---but do not establish national annual LDAR coverage
\citep{sngpl2024,ssgc2024}. Because the
2019 lower distribution factor requires both more than half plastic pipe and LDAR in use, the
public evidence does not resolve either national criterion sufficiently to select it.

OGRA reports \num{10.67} million tonnes of refinery product output. IPCC calls for crude feed,
so output converted at \SI{0.85}{\tonne\per\cubic\metre} is a provisional activity proxy with
\SI{20}{\percent} uncertainty. Produced gas is similarly used as provisional gas-processing
throughput. Neither proxy is hidden: both can be excluded by a coverage scenario.

Post-meter industrial activity is reconstructed directly from OGRA end-use controls rather than
one undifferentiated non-residential total. The included FY2023--24 classes are commercial,
general industry, fertilizer, captive power, cement and power stations, together
\num{2439.23} MMCF d$^{-1}$, which converts to \num{25211.03} million m$^3$ yr$^{-1}$;
the registered activity the leak factor rides is \num{25210.98} million m$^3$ yr$^{-1}$
before the fertilizer feedstock is withheld and \num{18684.62} after. The factor is defined
on gas consumed in combustion equipment, and \SI{83.5}{\percent} of the fertilizer class is
chemical feedstock, reformed and not burned, on the producers' metered feed and fuel gas
\citep{nfdc2025}; that feedstock, \num{6526.4} million m$^3$, is outside the base and the
process fuel stays in. Residential gas is
represented by the separate appliance fallback; CNG and unclassified other-system gas remain
outside the industrial/power row. The included class sum differs from the registered six-class
activity by only \num{0.042} million m$^3$ because published daily flows are rounded.

\begin{table*}[htbp]
\centering
\footnotesize
\begin{tabular}{@{}llrlp{6.0cm}@{}}
\toprule
\textbf{Segment} & \textbf{Activity} & \textbf{Value} & \textbf{Status} & \textbf{Principal structural issue} \\
\midrule
Oil production & onshore oil & 4,103.77 thousand m$^3$ & official & offshore split unavailable \\
Oil refining & converted product output & 12,552.94 thousand m$^3$ & proxy & crude-feed throughput unavailable \\
Gas production/gathering & produced gas & 32,309.47 million m$^3$ & official/proxy & well and field split unavailable \\
Gas processing & produced gas & 32,309.47 million m$^3$ & proxy & plant throughput and sour share \\
Transmission & pipeline length & 13,526.69 km & official & LDAR and seal practice \\
Distribution & pipeline length & 201,230 km & official & material share and LDAR \\
Post-meter appliances & domestic connections & 10,748,538 & official/proxy & one appliance per meter \\
Post-meter industry/power & non-residential gas & 25,210.98 million m$^3$ & official & factor transferability; fertilizer feedstock withheld from this base (see text) \\
\bottomrule
\end{tabular}
\caption{FY2023--24 reference activities. A proxy is estimated with an explicit fallback and
structural exclusion case; it is not described as observed.}
\label{tab:activity}
\end{table*}

\subsection{Observational and official comparison layers}

The current World Bank individual-flare workbook \citep{worldbankflaring2026}---built on
VIIRS Nightfire satellite pyrometry \citep{viirsnightfire}---contains annual volumes for
2012--2025 and \num{130} sites labelled Pakistan. The UNEP MARS snapshot contains plume and source tables under
a CC BY-NC-SA licence \citep{unepmars2026}; after filtering country and sector, it contains
\num{20} oil-and-gas sources and \num{174} plume detections. A public-domain geoBoundaries ADM1
layer assigns both products to provinces \citep{runfola2020}. Field-level evaluation uses
OGRA Table~3.34: its 84 published rows reconcile to \num{2637.75} MMCF d$^{-1}$, within
\num{0.03} MMCF d$^{-1}$ of the official total, and are canonicalized to 73 named fields.

The official comparison is extracted programmatically from Pakistan's 2022 CRT workbook. The
pipeline stores the package and workbook SHA-256 hashes, worksheet names, cell coordinates, raw
values, units, formulas and notation keys. The BTR's 2021 1.B.2 value is retained as a second
official-year comparator.

\subsection{Boundary and completeness}

The reference account covers oil and natural-gas systems under 1.B.2. Coal mining fugitives
(1.B.1) are outside the paper's stated boundary and are retained as NE in the completeness
matrix, now with a partial Tier-1 screen beside the key: the Thar mine-mouth deliveries the
power ledger holds bound the surface share (see
\texttt{outputs/fugitives\_coal\_screen\_2024.csv}), while underground production---whose
per-tonne factors run thirty times the surface low---remains unquantified. Oil and gas exploration, abandoned wells, natural-gas vehicles and anomalous events
also remain NE; storage is NO because the regulator describes facilities as planned rather
than operating. Venting and flaring are included elsewhere (IE) in the inclusive segment
factors and reported separately only as observation memos.
For every NE row, the completeness matrix records whether it is potentially material, probably
minor, unknown or outside scope and identifies the activity needed to estimate it. No numerical
upper bound is invented where public activity is unavailable.
